% 自编教学补充；阅读来源见本章 README。
\begin{frame}[fragile]{6.4 非金融小案例：左右走动}
\small 每步向左或向右移动一格，起点为 0。先生成步长，再用 cumsum 得到位置；随机种子使课堂输出可复现。
\medskip
\begin{lstlisting}
walk_rng = np.random.default_rng(7)
directions = walk_rng.integers(0, 2, size=12)
steps = np.where(directions == 0, -1, 1)
positions = np.r_[0, steps.cumsum()]
print(steps)
print(positions)
print("最远距离：", np.abs(positions).max())
\end{lstlisting}
\end{frame}
\begin{frame}[fragile]{课堂练习：6.4 非金融小案例：左右走动}
改成 100 步，并计算走动后位于起点右侧的比例；分母是否包括起点？
\medskip
\begin{block}{检查思路}
先写出预期结果，再运行。参考答案见配套 Notebook 的折叠单元格。
\end{block}
\end{frame}
